
\subsubsection{5.1 Framework Validation (RQ1)}

The measurement framework executes correctly end-to-end. All five modules (ast_decomposition, ast_metric, kl_metric, sep_analysis, statistical_tests) ran without errors on real DeepSeek-Coder-7B-instruct-v1.5 checkpoints.

\textbf{Proof-of-concept results (h-e1, synthetic data):}

The following results are derived from hand-crafted code completions with engineered structural differences and do not reflect real GRPO or DPO training outputs.

\begin{table}[htbp]
\centering
\begin{tabular}{lll}
\toprule
Metric & GRPO & DPO \\
\midrule
Mean semantic AST edit distance & 3.500 & 1.000 \\
Mean edit-per-KL (low-KL pairs) & ~25.9 & ~3.7 \\
Syntax pass rate & 6/6 (1.000) & 6/6 (1.000) \\
Bootstrap 95\% CI (differential) & [4.6500, 8.7314] & — \\
Mean differential & 6.5047 & — \\
Matched KL pairs (tolerance=0.15) & 27 & — \\
Bootstrap samples & 10,000 & — \\
\bottomrule
\end{tabular}
\end{table}

Per-problem semantic AST edit distances on the 6 proof-of-concept problems:

\begin{table}[htbp]
\centering
\begin{tabular}{lll}
\toprule
Problem & GRPO edit distance & DPO edit distance \\
\midrule
HumanEval/0 & 4.0 & 5.0 \\
HumanEval/1 & 6.0 & 0.0 \\
HumanEval/2 & 1.0 & 0.0 \\
HumanEval/3 & 0.0 & 0.0 \\
HumanEval/4 & 5.0 & 0.0 \\
HumanEval/5 & 5.0 & 1.0 \\
**Mean** & **3.500** & **1.000** \\
\bottomrule
\end{tabular}
\end{table}

The bootstrap 95\% CI [4.65, 8.73] excludes zero, confirming that the measurement infrastructure correctly detects a difference when one is engineered to exist. Per-pair edit-per-KL differentials ranged from 2.15 (pair 26, matching steps 1000/1000) to 22.18 (pair 2, matching steps 100/300).

Component-level verification:

\begin{table}[htbp]
\centering
\begin{tabular}{lll}
\toprule
Component & Status & Evidence \\
\midrule
FA-AST taxonomy (CF+DF classification) & Functional & SEP values in valid [0, 1] range \\
ZSS semantic edit distance & Functional & Correct distances on 6 problems \\
KL log matching (tolerance=0.15) & Functional & 27 matched pairs found \\
Bootstrap CI (10,000 samples) & Functional & CI excludes zero on PoC data \\
Mann-Whitney U test & Functional & Test executes; requires sufficient unique checkpoints \\
\bottomrule
\end{tabular}
\end{table}

\begin{figure}[htbp]
\centering
\includegraphics[width=\columnwidth]{figures/edit_per_kl_comparison.png}
\caption{Semantic-edit-per-KL comparison across KL-matched pairs (h-e1 proof-of-concept data)}
\label{fig:edit_per_kl_comparison}
\end{figure}

\textit{Figure 1. Semantic-edit-per-KL comparison between GRPO and DPO across 27 KL-matched pairs from h-e1 proof-of-concept (hand-crafted) data. GRPO exhibits higher raw AST semantic edit distance per unit KL divergence across all pairs.}

\begin{figure}[htbp]
\centering
\includegraphics[width=\columnwidth]{figures/bootstrap_ci_differential.png}
\caption{Bootstrap distribution of mean differential with 95\% CI}
\label{fig:bootstrap_ci_differential}
\end{figure}

\textit{Figure 2. Bootstrap distribution (10,000 samples) of the mean semantic-edit-per-KL differential. The 95\% CI [4.65, 8.73] excludes zero on the proof-of-concept data.}

\begin{figure}[htbp]
\centering
\includegraphics[width=\columnwidth]{figures/ast_edit_distances.png}
\caption{Per-problem AST semantic edit distances}
\label{fig:ast_edit_distances}
\end{figure}

\textit{Figure 3. Per-problem AST semantic edit distances for GRPO and DPO on six HumanEval+ problems (h-e1, synthetic data).}

\subsubsection{5.2 Preliminary SEP Analysis (RQ2)}

The following results use real checkpoints from h-m1 but are severely underpowered due to checkpoint aliasing (Section 5.3); they should not be interpreted as empirical evidence for or against the mechanism hypothesis.

\textbf{Table 1: Semantic Edit Proportion (SEP) — h-m1 Preliminary Analysis}

\begin{table}[htbp]
\centering
\begin{tabular}{llll}
\toprule
Condition & Mean SEP & N samples & vs. DPO \\
\midrule
GRPO-binary & 0.2371 & 192 & −0.0006 (lower) \\
GRPO-error-type & 0.2371 & 192 & −0.0006 (lower) \\
DPO & 0.2377 & 189 & — \\
\bottomrule
\end{tabular}
\end{table}

Mann-Whitney U (GRPO-binary vs. DPO): U = 18,346.5, p = 0.4248 (not significant at α = 0.05).
Effect size (GRPO − DPO): −0.0072.
Spearman correlation: undefined (NaN) due to zero variance in x-axis across aliased pairs.

The GRPO-binary and GRPO-error-type conditions produce identical SEP statistics because the h-m1 analysis reused checkpoints from h-e1; checkpoint aliasing caused both conditions to analyze the same checkpoint-100 files. The two reward functions are expected to produce distinguishable checkpoints only in a corrected run with dedicated full-scale training.

An observation of note is the divergence between the raw edit distance result and the SEP result. The proof-of-concept (h-e1) measured GRPO mean semantic edit distance of 3.5 versus DPO mean of 1.0 (a +250\% difference on synthetic data). The h-m1 analysis measured SEP ≈ 0.237 for both conditions. These two quantities measure different properties — raw edit volume versus the proportion of edits targeting semantic nodes — and derive from different experimental conditions (synthetic versus real checkpoints). They cannot be directly compared.

A dry-run smoke test on 10 fresh HumanEval+ problems (3 matched pairs) produced SEP ≈ 0.238 for both GRPO and DPO, providing weak corroborating evidence that the near-equality is not solely an artifact of aliasing, though this result is also underpowered.

\begin{figure}[htbp]
\centering
\includegraphics[width=\columnwidth]{figures/gate_sep_comparison.png}
\caption{SEP comparison between GRPO and DPO}
\label{fig:gate_sep_comparison}
\end{figure}

\textit{Figure 4. Semantic Edit Proportion (SEP) comparison between GRPO and DPO (h-m1, preliminary, underpowered). Distributions are nearly overlapping; medians are approximately 0.237 for both conditions.}

\begin{figure}[htbp]
\centering
\includegraphics[width=\columnwidth]{figures/ast_edit_distribution.png}
\caption{AST edit distance distribution by node category}
\label{fig:ast_edit_distribution}
\end{figure}

\textit{Figure 5. Distribution of AST edit distances across checkpoint pairs for GRPO and DPO, broken down by node category.}

\begin{figure}[htbp]
\centering
\includegraphics[width=\columnwidth]{figures/ast_node_heatmap.png}
\caption{AST node type heatmap}
\label{fig:ast_node_heatmap}
\end{figure}

\textit{Figure 6. Heatmap of AST node type frequencies modified by GRPO versus DPO across checkpoint steps.}

\subsubsection{5.3 Checkpoint Aliasing Confound (RQ3)}

The h-m1 analysis was compromised by a confound designated here as \textit{checkpoint aliasing}, in which insufficient checkpoint diversity caused a nominally large analysis to collapse to an effective sample size far smaller than the nominal count.

\textbf{What occurred:} The h-e1 proof-of-concept smoke test saved only 2 real GRPO checkpoints (step-100 and step-200). The h-m1 analysis was designed to use checkpoints from steps 100–1000 for its 27-pair KL-matched analysis. Checkpoint retrieval logic fell back to checkpoint-100 for steps 300–1000 when the requested files were not found.

\textbf{Consequence:} 25 of 27 analysis pairs aliased to checkpoint-100. The 192 GRPO SEP values were derived from effectively 2 distinct checkpoints, violating the independence assumption of the Mann-Whitney test. The nominal n = 192 corresponds to n_eff ≈ 2.

\textbf{Detection:} Figure 8 (sep_vs_kl_trajectory.png) shows a near-flat SEP trajectory across the 27 nominal pairs, inconsistent with output from a diverse 1000-step training run. Figure 7 (reward_correctness_scatter.png) shows tight clustering of reward values across checkpoint pairs. These patterns are characteristic of aliasing and would not appear in a correctly executed run with 10 or more unique checkpoints.

The confound is not detectable from pass@1 metrics alone. A team relying solely on benchmark outcomes would not detect this failure mode.

\textbf{Recommended safeguard:}

```python
unique_checkpoints = len(set(checkpoint_paths))
assert unique_checkpoints >= N_min, (
    f"Checkpoint aliasing detected: only {unique_checkpoints} unique "
    f"checkpoints found (required >= {N_min}). Abort analysis."
)
```

Recommended N_min = 10 for mechanism-level analysis, N_min = 5 for existence-level checks.

\begin{figure}[htbp]
\centering
\includegraphics[width=\columnwidth]{figures/reward_correctness_scatter.png}
\caption{Reward vs. correctness scatter}
\label{fig:reward_correctness_scatter}
\end{figure}

\textit{Figure 7. Scatter plot of execution reward signal versus pass@1 correctness for GRPO-trained checkpoints. Tight clustering of values is consistent with checkpoint aliasing.}

\begin{figure}[htbp]
\centering
\includegraphics[width=\columnwidth]{figures/sep_vs_kl_trajectory.png}
\caption{SEP vs. KL divergence trajectory}
\label{fig:sep_vs_kl_trajectory}
\end{figure}

\textit{Figure 8. SEP as a function of KL divergence trajectory across training steps. The near-flat line across 27 nominal pairs is attributable to checkpoint aliasing; only two unique GRPO checkpoints (step-100, step-200) were available.}
